%\begin{center}
% Está puesto en portada.tex
\thispagestyle{empty}
%forma de introducir imágenes. el \\[0.5 cm] de final de línea introduce un salto de ese tamaño.
%width=1\textwidth indica el tamaño de la imágen (valores entre 0-1). 
\includegraphics[width=1\textwidth]{figuras/cabecera.png}  \\[0.5 cm]

\LARGE UNIVERSIDAD POLITÉCNICA DE MADRID \\ [1 cm]

\LARGE ESCUELA TÉCNICA SUPERIOR DE INGENIERÍA Y DISEÑO INDUSTRIAL \\ [1 cm]

\LARGE Grado en Ingeniería Electrónica y Automática Industrial\\ [1 cm]

\LARGE \textbf{TRABAJO FIN DE GRADO}\\[1 cm]

\Huge \textsc{Integración de un brazo robótico en una red domótica}\\[2 cm]

\Large Firma Autor \\[3 cm]

%flushleft alinea a la izquierda el texto

\begin{multicols}{2} 
\begin{flushleft} 
\Large \emph{Firma Cotutor (si lo hay)}
\end{flushleft}

\begin{flushright} 
\Large \emph{Firma Tutor}
\end{flushright}

\end{multicols} 

%rellena de blanco el resto de la página para escribir abajo del todo
\vfill

\end{center}
\end{titlepage}


\cleardoublepage 